% chktex-file 44
% chktex-file 8
\documentclass[10pt,leqno]{amsart}
\usepackage{graphicx}
\baselineskip=16pt

\usepackage{indentfirst,csquotes}

\topmargin= .5cm
\textheight= 20cm
\textwidth= 32cc
\baselineskip=16pt

\evensidemargin= .9cm
\oddsidemargin= .9cm

\usepackage{amssymb,amsthm,amsmath}
\usepackage{xcolor,paralist,hyperref,titlesec,fancyhdr,etoolbox}
\newtheorem{theorem}{Theorem}[]
\newtheorem{definition}[theorem]{Definition}
\newtheorem{example}[theorem]{Example}
\newtheorem{lemma}[theorem]{Lemma}
\newtheorem{proposition}[theorem]{Proposition}
\newtheorem{corollary}[theorem]{Corollary}
\newtheorem{conjecture}[theorem]{Conjecture}

\titleformat{\section}[display]{\normalfont\huge\bfseries\centering}{\centering\chaptertitlename\thechapter}{10pt}{\Large}
\titlespacing*{\section}{0pt}{0ex}{0ex}

\hypersetup{ colorlinks=true, linkcolor=black, filecolor=black, urlcolor=black }
\def\proof{\noindent {\it Proof. $\, $}}
\def\endproof{\hfill $\Box$ \vskip 5 pt }

\usepackage{lipsum}

\begin{document}

\title{Twelve lonely runners}
\author[T. Sungkawichai]{Touch Sungkawichai}
\date{\today}
\email{touchsungkawichai@gmail.com}
\maketitle 

% \let\thefootnote\relax
% \footnotetext{MSC2020: Primary 00A05, Secondary 00A66.}

\begin{abstract}
  The Lonely Runner Conjecture asserts that for any number of runners moving at distinct constant speeds around a unit circle
  from a common starting point, each runner is lonely at some time. A runner is said to be lonely if its distance from every 
  other runner exceeds $1/n$, where $n$ is the total number of runners. 
  Tao\cite{Tao} showed that it suffices to verify the conjecture for finitely many positive integers. 
  Using computational method, Trakultongchai\cite{Tanupat} verified the conjecture for up to ten runners.
  In this work, I prove that the conjecture holds for twelve runners, assuming it holds for eleven runners.
\end{abstract}

\bigskip

\section*{Introduction}

\begin{thebibliography}{99}

\bibitem{Tao} Spiegel, M. R. (1981). Theory and problems of Advanced Calculus: Si (metric) edition. McGraw-Hill. 
\bibitem{Tanupat} Spiegel, M. R. (1981). Theory and problems of Advanced Calculus: Si (metric) edition. McGraw-Hill. 

\end{thebibliography}
\end{document}



% Write $S_k$ as the set modulo $P$ such that $C_k[t, v] = 1 \iff t \cdot v \in S_k$ where $C_k$ is the coverage bitmask. 
%  
% \begin{lemma}
% For any solution $\set{v_1, \dots, v_k}$, there exists a solution $\set{1, u_2, \dots, u_k}$
% and a scalar $u \in (\mathbb{Z}/P\mathbb{Z})^\times$ such that
% \[
%     \set{v_1, \dots, v_k} = \set{u,\, u u_2,\, \dots,\, u u_k}.
% \]
% \end{lemma}
%  
% \begin{proof}
% Set $u = v_1$ and $u_i = v_1^{-1} v_i \pmod{P}$ for each $i$, so that $u_1 = 1$ and
% $u u_i = v_1 \cdot v_1^{-1} v_i = v_i$. Hence
% \[
%     \set{u, u u_2, \dots, u u_k} = \set{v_1, v_2, \dots, v_k}.
% \]
% It remains to show that $\set{1, u_2, \dots, u_k}$ is a valid solution.
% Fix any $t \in (\mathbb{Z}/P\mathbb{Z})^\times$ and set $t' = t v_1^{-1}$.
% Since $\set{v_1, \dots, v_k}$ is a solution, there exists some index $i$ such that
% \[
%     t' \cdot v_i \in S_k.
% \]
% But $t' \cdot v_i = t v_1^{-1} v_i = t \cdot u_i$, so $t \cdot u_i \in S_k$,
% i.e.\ $C_k[t, u_i] = 1$. Since $t$ was arbitrary, $\set{1, u_2, \dots, u_k}$ is a valid solution.
% \end{proof}
